% F4: um passe real atravessando as camadas (o passe de Mac Allister no gol de Di María).
% Fonte: docs/02-dos-dados-ao-grafo.md (trechos do JSON, da camada 0 e do grafo). Nomes abreviados na tabela.
\providecommand{\jq}{\char34{}}
\providecommand{\ji}{\hspace*{1.2em}}
\vspace{2mm}
\noindent\begin{minipage}{\columnwidth}
\scriptsize
\textit{(a) Evento no JSON bruto}\par\smallskip
{\ttfamily\frenchspacing
\renewcommand{\arraystretch}{1.0}%
\begin{tabular}{@{}l@{}}
\{\jq minute\jq: 35, \jq second\jq: 20, \\
\ji\jq type\jq: \{\jq name\jq: \jq Pass\jq\}, \\
\ji\jq player\jq: \{\jq name\jq: \jq Alexis Mac Allister\jq\}, \\
\ji\jq pass\jq: \{\jq recipient\jq: \\
\ji\{\jq name\jq: \jq Ángel Fabián Di María Hernández\jq\}, \\
\ji\ji\jq goal\_assist\jq: true\}\}
\end{tabular}}\par\medskip
\textit{(b) Linhas da tabela de ações}\par\smallskip
\setlength{\tabcolsep}{2pt}%
\begin{tabular}{@{}llllllll@{}}
\toprule
min. & jogador & ação & certo & zonas & xT & recebedor & fase \\
\midrule
36 & Mac Allister & passe & sim & 74$\to$92 & 0,184 & Di María & 132 \\
36 & Di María & finalização & gol & 92$\to$91 & 0 & -- & 132 \\
\bottomrule
\end{tabular}\par\medskip
\textit{(c) Aresta no grafo}\par\smallskip
\centering
\begin{tikzpicture}[box/.style={draw, rounded corners=2pt, inner sep=3pt}]
\node[box] (a) at (0,0) {Mac Allister};
\node[box] (b) at (0.62\columnwidth,0) {Di María};
\draw[-{Stealth[length=2mm]}] (a) -- node[above] {\texttt{PASSOU\_PARA}} node[below, align=center] {minuto 36, zonas 74$\to$92\\ xT 0,184, fase 132} (b);
\end{tikzpicture}
\captionof{figure}{Um passe real nas três camadas: o de Mac Allister para o gol de Di María, de evento no JSON a linha da tabela e aresta no grafo. O minuto do JSON conta minutos decorridos; o da tabela, o da transmissão.}\label{fig:lance}
\end{minipage}\par\smallskip
